% Boundary-flow extension of the mathematical forward model.

\section{Boundary Flows and Separate Boarding/Alighting Processes}
\label{sec:boundary-flow-observation-processes}

The ordinary OD--time model assumes that every passenger contributing to an
observation enters the modelled window through one of its OD cells.  That
assumption is not always appropriate when the window is a segment of a
service, rather than a complete day.  Some passengers can already be on board
when the window begins, while others can board at its end and alight later.
The boundary-flow extension represents these two populations explicitly while
leaving the ordinary OD demand model intact.  It also permits boarding and
alighting rows to have different observation scales.

\subsection{Additive forward equation}

Let $j=1,\ldots,n$ index the free OD--time cells and let
$x_j(\vartheta)\geq 0$ be the demand generated by the selected gravity model
with parameter vector $\vartheta$.  Let $g=1,\ldots,n_I$ index groups of
passengers already on board at the beginning of the observation window, and
let $h=1,\ldots,n_T$ index groups boarding at the terminal boundary.  Their
non-negative latent flows are denoted by
$z^{I}_g(\alpha)$ and $z^{T}_h(\gamma)$, respectively.

The primary measurement rows are kept in one canonical order.  The ordinary
OD response, the two boundary responses, and the known fixed contribution are
represented by
\[
 A\in\mathbb{R}^{m\times n},\qquad
 B_I\in\mathbb{R}^{m\times n_I},\qquad
 B_T\in\mathbb{R}^{m\times n_T},\qquad
 \bm c\in\mathbb{R}^{m}.
\]
The latent expected measurement vector is
\begin{equation}
  \bm\ell(\vartheta,\alpha,\gamma)
   = A\bm x(\vartheta)
     +B_I\bm z^{I}(\alpha)
     +B_T\bm z^{T}(\gamma)
     +\bm c.
  \label{eq:boundary-additive-latent-measurements}
\end{equation}

The columns of $B_I$ and $B_T$ are not OD cells.  They are case-defined
boundary cohorts, and their entries describe the expected contribution of one
cohort to every measurement row.  Consequently, an initial-onboard cohort
may contribute to alighting rows at any later stop; it is not restricted to
the first alighting row.  Similarly, a terminal-outflow cohort contributes to
boarding rows at the end of the window and is not interpreted as ordinary
OD demand inside the window.

When no boundary blocks are present, equation~\eqref{eq:boundary-additive-latent-measurements}
reduces exactly to the existing OD-only expression
$\bm\ell=A\bm x+\bm c$.

\subsection{Separate observation processes}

For each row $r$, let $q(r)\in\{\mathrm{boarding},\mathrm{alighting}\}$
denote its event type.  The current observation model uses one global scale
$\rho>0$ and optional fixed type-specific scales
\[
 s_{\mathrm{boarding}}>0,\qquad
 s_{\mathrm{alighting}}>0.
\]
The unfloored mean is therefore
\begin{equation}
  \widetilde\mu_r
   =\rho\,s_{q(r)}\,\ell_r.
  \label{eq:boundary-observation-mean}
\end{equation}
For numerical safety the implemented mean is floored at a declared positive
value $\mu_{\min}$:
\[
 \mu_r=\begin{cases}
   \widetilde\mu_r,&\widetilde\mu_r>\mu_{\min},\\
   \mu_{\min},&\widetilde\mu_r\leq\mu_{\min}.
 \end{cases}
\]
Rows on the floor have a constant mean and therefore contribute no derivative
through the measurement response.  The scales are fixed inputs in the
current implementation; they are not additional fitted parameters.

The count likelihood is applied to the resulting means, for example
\begin{align}
  y_r\mid\mu_r &\sim \operatorname{Poisson}(\mu_r),
  &&\text{or} &
  y_r\mid\mu_r,\phi &\sim
       \operatorname{NegBinomial}(\mu_r,\phi),
\end{align}
with the declared calibration mask determining which rows enter the
likelihood.  A row can be supported by the OD response, by one or both
boundary responses, or by a known fixed contribution.  A positive row with no
active support is not an observed zero: it is a support or data-contract
failure unless an explicit calibration policy excludes it.

\subsection{Latent-flow parameterizations}

The boundary flows need not have one free parameter for every vehicle and
stop.  A direct grouped parameterization uses one raw parameter per declared
cohort:
\begin{equation}
 z^{I}_g(\alpha_g)=\epsilon+\operatorname{softplus}(\alpha_g),
 \qquad
 z^{T}_h(\gamma_h)=\epsilon+\operatorname{softplus}(\gamma_h),
 \label{eq:boundary-direct-positive-flow}
\end{equation}
where $\epsilon>0$ is a numerical positivity floor.  A lower-dimensional
parameterization uses a fixed basis $P_I$ or $P_T$ and applies the same
non-negative transformation to the generated flows, for example
\[
 \bm z^{I}=\epsilon\bm 1+\operatorname{softplus}(P_I\bm\alpha),
 \qquad
 \bm z^{T}=\epsilon\bm 1+\operatorname{softplus}(P_T\bm\gamma).
\]
The basis encodes scientific structure such as route, direction, time band,
or a small number of survival-shape factors.  It is not a replacement for
the boundary operators: $B_I$ and $B_T$ still determine how those flows are
seen by the detectors.

The optimizer sees one deterministic vector
\[
 \bm\eta=(\vartheta,\alpha,\gamma),
\]
with stable names and slices for the OD, initial-onboard, and terminal-outflow
blocks.  If a block carries a quadratic penalty, its contribution can be
written as
\[
 \mathcal R_I(\alpha)=\tfrac12\lambda_I\|\alpha\|^2,
 \qquad
 \mathcal R_T(\gamma)=\tfrac12\lambda_T\|\gamma\|^2,
 \quad \lambda_I,\lambda_T\geq0.
\]
The complete MAP objective is then
\begin{equation}
  \mathcal J(\bm\eta)
   =-\ell_{\mathrm{count}}(\bm\eta)
     +\mathcal R_{\mathrm{OD}}(\vartheta)
     +\mathcal R_I(\alpha)+\mathcal R_T(\gamma).
  \label{eq:boundary-map-objective}
\end{equation}
For ML, all declared prior or regularization terms are omitted.

\subsection{Gradient structure and identifiability}

Let $\bm w=\partial\mathcal J/\partial\bm\ell$ be the row-space derivative
of the count objective.  The derivative with respect to the latent initial
flow is obtained through the transpose product
\[
 \frac{\partial\mathcal J}{\partial\bm z^{I}}=B_I^{\mathsf T}\bm w,
 \qquad
 \frac{\partial\mathcal J}{\partial\bm z^{T}}=B_T^{\mathsf T}\bm w,
\]
followed by the Jacobian of the chosen flow parameterization.  This is the
same forward/adjoint linear-operator pattern as for the OD response.  The
boundary blocks therefore add parameter directions to the statistical model,
but do not require a second route-choice calculation.

Unrestricted trip-by-stop boundary flows are generally underidentified: many
combinations of initial stocks and alighting hazards can produce the same
aggregate stop counts.  Groups or low-dimensional bases must consequently be
chosen from a declared observation window and supported by regularization or
other prior information.  The case adapter owns the interpretation of a
``trip start'', the boundary cohorts, and the row-to-column construction; the
public mathematical contract does not infer those meanings from stop names.

The two event scales are likewise not freely identifiable when a global
production or detection scale is estimated.  The current model avoids that
confounding by treating boarding and alighting scales as fixed.  An eventual
estimated-scale extension must include an explicit identifying constraint,
such as fixing one event scale or constraining the product of the two scales.

\subsection{Relation to auxiliary journey observations}

Boundary blocks change the primary mean in equation~\eqref{eq:boundary-observation-mean}.
They are therefore different from an auxiliary journey-attribute observation
term (for example, a GPS-derived travel-time histogram).  An auxiliary term
adds a separate likelihood contribution based on the same demand vector; it
does not alter the primary boarding/alighting mean.  The two mechanisms may
be combined, but their contributions, masks, and provenance must remain
separate.

\subsection{Identity and reuse of response operators}

The OD response $A$ and each boundary response $B_I,B_T$ have separate
mathematical identities.  If the timetable, route choices, canonical OD
index, measurement-row order, and numeric type are unchanged, an existing
OD response can be reused when boundary blocks are added.  A boundary block
is compatible only when its row order, column groups, matrix content, and
parameter/model specification match the fit.  Changing a boundary block or
its grouping changes the joint model identity and requires a new fit or
checkpoint, but it does not by itself require rebuilding $A$.

The current public implementation provides generic additive blocks, fixed
boarding/alighting scales, deterministic joint layouts, independent boundary
artifact identities, and contribution-level reporting.  Estimated event
scales, type-specific dispersions, and richer route-wide survival or hazard
models remain future extensions; they must preserve the identifiability and
provenance conditions stated above.
